\documentclass{../latex/utmreport}
\newcommand{\ReportTitle}{Lab 7: Three Small Programs: Processes, Redirection, a Pipe}
% Shared title-page and PDF metadata for every laboratory report.
\newcommand{\StudentName}{Gafenco Victor}
\newcommand{\StudentGroup}{FAF-241}
\newcommand{\CourseName}{Operating Systems}
\newcommand{\InstructorName}{Patricia Reițman}
\newcommand{\InstructorTitle}{Assistant lecturer}
\newcommand{\ReportCity}{Chișinău}
\newcommand{\ReportYear}{2026}
\newcommand{\CourseRepository}{https://github.com/yoda-dinmd/os-labs}

\selectlanguage{english}
\utmlogo{../latex/assets/UTM_Logo_English.jpeg}
\institution{Technical University of Moldova}
\faculty{Faculty of Computers, Informatics, and Microelectronics}
\department{Department of Software Engineering}
\studyprogramme{Software Engineering Study Programme}
\documenttype{Laboratory report}
\student{\StudentName}{\StudentGroup}
\supervisor{\InstructorName}{\InstructorTitle}
\renewcommand{\UTMCoordinatorLabel}{Instructor}
\cityyear{\ReportCity}{\ReportYear}
\reportsubtitle{\CourseName}
\hypersetup{pdfauthor={\StudentName},pdfsubject={\CourseName},
  pdftitle={\ReportTitle}}
\reporttitle{\ReportTitle}
\addto\captionsenglish{\renewcommand{\figurename}{Figure}\renewcommand{\tablename}{Table}\renewcommand{\bibname}{References}}
\renewcommand{\figurename}{Figure}
\renewcommand{\tablename}{Table}
\renewcommand{\bibname}{References}
\usepackage{float}
\usepackage{placeins}
\setlength{\emergencystretch}{2em}
\newcommand{\code}[1]{\texttt{\detokenize{#1}}}
% Keep screenshots beside their observations, with bounded height and width.
\newcommand{\screenshot}[4][\textwidth]{%
  \begin{figure}[H]
    \centering
    \includegraphics[width=#1,height=0.62\textheight,keepaspectratio]{screenshots/#2}
    \caption{#3}\label{#4}
  \end{figure}}
\newcommand{\reportreferences}[2]{%
  \clearpage\phantomsection\addcontentsline{toc}{chapter}{References}
  \begin{thebibliography}{9}
    \bibitem{handout} Technical University of Moldova, FCIM / FAF.
    \textit{#1}. Operating Systems laboratory handout, academic year 2026--2027.
    Supplied course document: \code{#2}.
  \end{thebibliography}}

\providecommand{\tightlist}{\setlength{\itemsep}{0pt}\setlength{\parskip}{0pt}}
\begin{document}
\makeutmreporttitlepage
\tableofcontents
\clearpage
\begin{utmabstract}[Abstract]
Four Linux programs demonstrate process creation, program replacement, delayed child reaping, standard-output redirection and a two-stage pipeline. Successful and failed execution, a visible zombie and a deliberately retained pipe writer connect the observed behaviour to the underlying process and descriptor operations.
\utmkeywords{Linux, fork, exec, waitpid, zombie, redirection, pipe}
\end{utmabstract}
\utmintroduction

Four C programs examine process creation, program replacement, child
reaping, output redirection and communication through a pipe on Linux.
The submitted \texttt{run.c} runs \texttt{ls\ -l}; the \texttt{date} and
nonexistent-program versions were also compiled and executed.
\texttt{zombie.c} delays reaping by 20 seconds, \texttt{redirect.c}
sends a listing to a file, and \texttt{pipeline.c} connects
\texttt{ls\ /etc} to \texttt{wc\ -l}.

All four programs compiled without diagnostics using GCC (SUSE Linux)
16.2.1 with \texttt{-Wall\ -Wextra\ -Werror}. Each can be built
independently, for example:

\begin{lstlisting}
gcc -Wall -Wextra -o run run.c
./run
\end{lstlisting}

The complete command outputs are in \texttt{evidence/}. Process
identifiers, directory contents and timestamps describe these runs and
can differ on another machine. The experiments assume standard
descriptors 0, 1 and 2 are initially open.

\chapter{Process creation and program
replacement}\label{process-creation-and-program-replacement}

\section{The three steps}\label{the-three-steps}

\texttt{fork()} creates a child and returns its PID to the parent, zero
to the child, or a negative value on failure. Both processes continue
after the call. In the child, \texttt{execlp()} searches for an
executable using \texttt{PATH} and replaces the current program image.
The process retains its PID. The parent calls
\texttt{waitpid(pid,\ \&status,\ 0)} to collect this particular child's
termination status.

The initial \texttt{fflush(stdout)} prevents buffered parent output from
being copied and printed twice after \texttt{fork}. The child's flush
preserves its diagnostic before \texttt{exec} replaces its address
space. The final parent message follows the child's termination because
of \texttt{waitpid}.

\section{Observed successful and failed
execution}\label{observed-successful-and-failed-execution}

Selected output from \texttt{evidence/run.txt}:

\begin{lstlisting}
$ ./run
parent: my PID is 27
child:  my PID is 28, my parent is 27
total 108
[listing omitted]
parent: child 28 ended with status 0
\end{lstlisting}

Replacing the execution line with
\texttt{execlp("date",\ "date",\ (char\ *)NULL);} and recompiling
produced:

\begin{lstlisting}
parent: my PID is 29
child:  my PID is 35, my parent is 29
Mon Oct 26 12:01:04 EET 2026
parent: child 35 ended with status 0
\end{lstlisting}

Replacing it with
\texttt{execlp("nosuchprog",\ "nosuchprog",\ (char\ *)NULL);} produced:

\begin{lstlisting}
parent: my PID is 36
child:  my PID is 42, my parent is 36
exec: No such file or directory
parent: child 42 ended with status 127
\end{lstlisting}

\begin{enumerate}
\def\labelenumi{\arabic{enumi}.}
\tightlist
\item
  \texttt{perror("exec")} runs only after a failed \texttt{execlp}:
  successful execution never returns to the old program. The nonexistent
  executable causes the call to return with an error.
\item
  The child printed 27 as ``my parent'' in the first run. This is the
  original process's PID, returned in the child by \texttt{getppid()},
  and it matches the parent's first message.
\item
  The parent observed status 0 for \texttt{ls} and 127 for the failed
  execution. The latter comes from the explicit \texttt{exit(127)} after
  \texttt{perror}, not directly from \texttt{execlp}. The \texttt{date}
  case also returned 0.
\end{enumerate}

Linux encodes the wait status. These children exited normally, so \texttt{WEXITSTATUS(status)} extracts their exit codes; signal termination would require separate handling.

\chapter{A zombie process}\label{a-zombie-process}

\texttt{zombie.c} is the original \texttt{ls} version with
\texttt{sleep(20);} immediately before \texttt{waitpid}. During this
interval, \texttt{ls} finishes but the parent has not yet collected its
status. The process table was filtered to the two relevant PIDs:

\begin{lstlisting}
$ gcc -Wall -Wextra -o zombie zombie.c
$ ./zombie > /dev/null &
$ sleep 2
$ ps -o pid,ppid,stat,cmd -p 43,44
    PID    PPID STAT CMD
     43       2 S    ./zombie
     44      43 Z    [ls] <defunct>
\end{lstlisting}

After the parent's 20-second sleep and subsequent wait:

\begin{lstlisting}
$ ps -o pid,ppid,stat,cmd -p 43,44
    PID    PPID STAT CMD
\end{lstlisting}

\begin{enumerate}
\def\labelenumi{\arabic{enumi}.}
\tightlist
\item
  \texttt{ls} has state \texttt{Z}, meaning zombie: it has terminated,
  but its exit information remains available for the parent. It is no
  longer executing instructions. The parent is sleeping (\texttt{S}).
\item
  The \texttt{{[}ls{]}\ \textless{}defunct\textgreater{}} entry is
  absent after the wait. The line \texttt{waitpid(pid,\ \&status,\ 0);}
  reaps the child. \texttt{sleep(20)} only delays that operation; it
  does not reap anything. The parent also exits after printing its final
  message.
\end{enumerate}

\chapter{Output redirection}\label{output-redirection}

In the child, \texttt{close(1)} releases standard output.
\texttt{open("out.txt",\ O\_WRONLY\ \textbar{}\ O\_CREAT\ \textbar{}\ O\_TRUNC,\ 0644)}
opens or creates the destination and truncates any previous contents.
Its requested creation permissions are filtered by the process's umask.
The diagnostic uses standard error, so it remains visible on the
terminal.

\begin{lstlisting}
$ gcc -Wall -Wextra -o redirect redirect.c
$ ./redirect
open gave descriptor 1
done: the listing is in out.txt
$ head -3 out.txt
total 108
-rw-r--r--. 1 victor victor     0 Oct 26 12:01 out.txt
-rwxr-xr-x. 1 victor victor 20456 Oct 26 12:01 pipeline
\end{lstlisting}

\begin{enumerate}
\def\labelenumi{\arabic{enumi}.}
\tightlist
\item
  \texttt{open} returned 1 because it allocates the lowest unused
  descriptor, and the child had just closed 1 while 0 and 2 remained
  open.
\item
  \texttt{ls} still writes to descriptor 1. Open descriptors survive
  \texttt{exec} unless marked close-on-exec, so descriptor 1 now
  identifies \texttt{out.txt}. The parent's separate descriptor table
  still directs its ``done'' message to the terminal.
\item
  To redirect errors with the same technique, close descriptor 2 before
  opening \texttt{err.txt}; the newly opened file then becomes standard
  error.
\end{enumerate}

The zero size shown for \texttt{out.txt} within its own listing reflects
the metadata read by \texttt{ls} before its buffered listing was
written. Reading the file afterward confirms that it contains output.

\chapter{A two-process pipeline}\label{a-two-process-pipeline}

\section{Descriptor ownership and
result}\label{descriptor-ownership-and-result}

\texttt{pipe(p)} supplies a read end \texttt{p{[}0{]}} and a write end
\texttt{p{[}1{]}}. Both children inherit them. The \texttt{ls} child
closes descriptor 1 and calls \texttt{dup(p{[}1{]})}, making standard
output refer to the pipe. The \texttt{wc} child closes descriptor 0 and
calls \texttt{dup(p{[}0{]})}, making standard input refer to the pipe.
Both children close the original pipe descriptors before \texttt{exec}.
The parent closes both ends and waits twice.

\begin{lstlisting}
gcc -Wall -Wextra -o pipeline pipeline.c
$ ./pipeline
213
$ ls /etc | wc -l
213
\end{lstlisting}

The two numbers were equal: both commands counted 213 output lines from
\texttt{ls\ /etc} in the test environment.

\section{Keeping a writer open}\label{keeping-a-writer-open}

The parent's \texttt{close(p{[}1{]});} was commented out in a separate
experimental version and the program was recompiled. It produced no
output and remained running. The process table after two seconds showed:

\begin{lstlisting}
$ ./pipeline-broken
(no output; still running after 2 seconds)
$ ps -o pid,ppid,stat,cmd -p 61,63
    PID    PPID STAT CMD
     61       2 Ss   ./pipeline-broken
     63      61 S    wc -l
\end{lstlisting}

The experiment was stopped with SIGINT, the signal produced by Ctrl-C.
The submitted \texttt{pipeline.c} retains the correct close.

\begin{enumerate}
\def\labelenumi{\arabic{enumi}.}
\tightlist
\item
  The working program and the shell both printed 213.
\item
  \texttt{wc} is the child that never finishes. \texttt{ls} has already
  exited, while the parent is waiting for \texttt{wc}.
\item
  \texttt{wc} is waiting for end-of-file, which cannot arrive while the
  parent still holds an open reference to the pipe's write end.
\end{enumerate}

An empty pipe is not necessarily at EOF: a blocking read waits if any
writer remains. Closing unused descriptors is therefore part of the
communication protocol, not merely housekeeping.

\utmconclusions

The experiments demonstrate how a shell combines process and descriptor
operations. \texttt{fork} creates a child, \texttt{exec} replaces its
program, and \texttt{waitpid} collects its result. Delaying the wait
exposes a zombie. Redirection works by changing the child's descriptor
table before execution, and a pipeline connects two otherwise unchanged
programs. The stalled-pipe experiment confirms that EOF depends on
closing every reference to the writing end.

\end{document}
